\chapter{Introduction}
\label{ch:intro}

\section{Motivation}

Legged robots move well on land, and many of the places where they are useful, such as shorelines, flooded
areas, ponds and river banks, also contain open water. A quadruped that crosses such water needs some way to
propel itself when its feet no longer touch the ground. Adding thrusters is one option \cite{Yao2023}, but
thrusters add mass, volume and actuators that are useless on land. The alternative is to swim with the legs
themselves \cite{Li2019,Qu2025,Han2025,Wang2025}. Then no extra hardware is needed, and swimming becomes a
question of how the existing joints are moved and what the feet look like.

Animals that swim with legs use two different physical principles \cite{Fish1996}. A dog, a muskrat or a duck
\emph{paddles}. It pushes water backwards with a foot oriented normal to its motion and then brings the foot
forward again \cite{Fish1984,Fish2021DogPaddle}. The thrust comes from drag on the foot during the power
stroke, and the recovery stroke unavoidably costs some thrust. Dolphins, sea turtles and many fish instead
\emph{flap} a fin or fluke across the swimming direction. The fin meets the oncoming water at a small angle of
attack and produces lift, whose forward component is the thrust. Drag-based paddling is simple, works at rest
and accelerates well. Lift-based flapping is two to three times more efficient at cruising speed
\cite{Fish1996,Walker2000}. Mammals that became fully aquatic changed over from the first principle to the
second.

Most swimming legged robots use the first principle, because a walking foot is already a paddle. This thesis
studies whether that choice is right for a small amphibious quadruped, BODY2, and how the leg should move for
each principle. The robot's legs were designed for walking. Each is a planar two-degree-of-freedom linkage
driven by two hobby servos, and the foot is a fan-shaped paddle that folds during the recovery stroke. Its
current swimming gait was taken from video of the robot and has not been designed for thrust. Whether this
paddle can be made efficient, or whether a small lift-producing fluke on the same leg would be better, cannot
be decided from the literature alone, because the answer depends on the leg geometry, the actuator limits and
the swimming speed.

\section{Problem Statement and Scope}

The problem is to find, for a given leg and given actuators, the propulsor and the periodic leg motion that
propel the robot most effectively, and to quantify how the answer depends on swimming speed. Three
properties make the problem difficult.

First, the flow is unsteady and separated. A paddle starts and stops twice per cycle, sheds a vortex at every
stop, and meets its own wake. A flapping fluke works at Reynolds numbers of $10^3$--$10^4$, where
leading-edge vortices and the phase between heave and pitch decide the thrust
\cite{Triantafyllou1991,Anderson1998}. Quasi-steady models capture trends but not the magnitudes of these
effects \cite{Kim2011}.

Second, the answer depends on speed. At rest, a paddle pushes water that is not moving and produces thrust
over most of its power stroke. At speed, the oncoming flow reduces the relative velocity during the power
stroke and increases it during the recovery. A lift-based foil behaves in the opposite way: its angle of
attack is set by the ratio of heave velocity to swimming speed, so it depends on speed through the kinematics.
A comparison at a single speed can therefore favour either principle.

Third, no single model is both accurate and cheap. A three-dimensional CFD simulation of one leg takes hours to
days and cannot be placed inside an optimizer that needs thousands of evaluations. A reduced-order model can,
but its coefficients are not calibrated and its lift model is a textbook flat plate.

Three simulation tools are therefore used, each for a different question, and none is expected to reproduce the
numbers of the others:
\begin{itemize}
  \item a \emph{reduced-order} whole-robot model in MuJoCo with CMA-ES gait optimization answers a question of
        \emph{mechanism}: which propulsion principle does an optimizer select under actuator and stability constraints? Its
        coefficients are not calibrated, so it is used for direction, not for magnitude;
  \item \emph{single-leg CFD} of the real BODY2 leg geometry in OpenFOAM with overset grids answers the questions of
        \emph{magnitude}: how much thrust and power each propulsor produces across forward speed, and where one overtakes the
        other. It is validated against an experiment and provides every quantitative result of this thesis;
  \item the \emph{whole-robot} real-time solver Flare Live provides the reference fluke motion. Its fluke force is checked against
        the CFD and found too low, so its results are used for trends only.
\end{itemize}
\Cref{fig:overview} shows how the three tools and their results are connected. Where the reduced-order model and the CFD
disagree, the CFD is authoritative; the disagreement itself is analysed in \cref{sec:disc-models}.

\begin{figure}[t]
  \centering
  \begin{tikzpicture}[>=Latex,font=\small,
      box/.style={draw=TUMblue,thick,rounded corners=2pt,align=center,inner sep=4pt,text width=35mm,minimum height=17mm,fill=TUMblue!6,execute at begin node={\hyphenpenalty=10000\relax}},
      res/.style={draw=black!45,rounded corners=2pt,align=center,inner sep=3pt,text width=35mm,minimum height=9mm,fill=black!3,font=\footnotesize,execute at begin node={\hyphenpenalty=10000\relax}},
      lab/.style={font=\footnotesize\itshape,text=black!60}]
    \node[box,font=\small] (mj) at (0,0) {\textbf{Reduced-order model}\\MuJoCo, Morison\\forces, servo model,\\CMA-ES (\cref{ch:mujoco})};
    \node[box] (cfd) at (5.1,0) {\textbf{Single-leg CFD}\\OpenFOAM, overset grids,\\real leg geometry\\(\cref{ch:cfd,ch:verification})};
    \node[box] (fl) at (10.2,0) {\textbf{Whole-robot model}\\Flare Live, lattice\\Boltzmann, real time\\(\cref{sec:cfd-flare})};
    \node[lab,above=1pt of mj.north] {fast, uncalibrated};
    \node[lab,above=1pt of cfd.north] {resolved, validated};
    \node[lab,above=1pt of fl.north] {whole robot, coarse};
    \node[res] (r1) at (0,-2.3) {mechanism: which principle\\an optimizer selects (RQ1)};
    \node[res] (r2) at (5.1,-2.3) {thrust, power, efficiency\\vs.\ speed; crossover (RQ2)};
    \node[res] (r3) at (10.2,-2.3) {cross-check of the fluke force (\cref{sec:ver-flare})};
    \node[res] (sp) at (5.1,-4.2) {self-propulsion estimate $4\bar T(U)=D(U)$ (RQ3)};
    \node[res] (hw) at (10.2,-4.2) {first hardware test\\(\cref{sec:res-hw})};
    \draw[->] (mj) -- (r1); \draw[->] (cfd) -- (r2); \draw[->] (fl) -- (r3); \draw[->] (r2) -- (sp);
    \draw[->,dashed] (mj.east) -- node[above,lab,fill=white,inner sep=1pt]{lift?} (cfd.west);
    \draw[->,dashed] (fl.west) -- node[above,lab,fill=white,inner sep=1pt]{L0 motion} (cfd.east);
    \draw[->,dashed] (r2.east) -- node[above,lab,fill=white,inner sep=1pt]{fluke force} (r3.west);
    \draw[<->,dashed] (sp.east) -- node[above,lab,fill=white,inner sep=1pt]{compare} (hw.west);
  \end{tikzpicture}
  \caption[Overview of the method]{Overview of the three simulation tools and their roles. Solid arrows: results; dashed arrows:
    information passed between tools. The reduced-order model motivates the lift-type propulsor, Flare Live provides the reference fluke motion, and the CFD
    provides the thrust curves for the self-propulsion estimate and the reference against which the Flare Live fluke force is checked.}
  \label{fig:overview}
\end{figure}

The scope is limited as follows.
\begin{itemize}
  \item \emph{Environment.} The robot swims at the surface of still water. Waves, currents and the transition
        between land and water are excluded. All propulsor CFD is single-phase with a rigid lid; the free surface is
        not modelled.
  \item \emph{Motion.} Only steady, periodic, straight-ahead swimming is considered. Turning, starting and
        feedback control are excluded.
  \item \emph{Propulsors.} Two propulsors fitted to the same leg are compared: the robot's folding fan paddle
        (drag-based) and a small swept fluke on a pin joint (lift-based). The fluke pitch is prescribed in
        the simulations; the passive pin-and-spring realization of the hardware is not simulated.
  \item \emph{Evaluation.} Thrust and hydrodynamic input power are computed for the propulsor alone (fan or
        fluke), without the linkage. Swimming speeds are predicted from a quasi-steady force balance with a
        range of hull resistances. Hardware experiments are limited to a first timed test from rest.
\end{itemize}

\section{Research Questions and Hypotheses}

The central question of this thesis is:
\begin{quote}
\itshape Which propulsion principle, drag-based paddling or lift-based flapping, lets the legs of BODY2 swim
faster and more efficiently, and how should each propulsor move?
\end{quote}
It is split into three research questions, each paired with a hypothesis.

\begin{description}[leftmargin=12mm,style=nextline]
\item[(RQ1)] \emph{Reduced-order optimization.} When a whole-robot gait is optimized in a reduced-order
  hydrodynamic model under actuator-speed and attitude limits, which propulsion mechanism does the optimizer
  select, and does the answer depend on the physics included in the model?

  \emph{Hypothesis H1.} If the model contains a lift term, the optimizer uses lift-based strokes, and at equal
  speed these need less power than drag-based strokes.

\item[(RQ2)] \emph{Drag versus lift across speed.} For the real leg geometry, how large are thrust, power and
  Froude efficiency of an optimized drag-type paddle and of a lift-type fluke between rest and
  \SI{0.4}{\metre\per\second}, and what limits each propulsor?

  \emph{Hypothesis H2.} The optimized paddle produces more thrust at low speed, but its thrust falls steeply
  with speed because the recovery stroke meets an increasing oncoming flow. The fluke produces less thrust at
  rest, keeps a larger fraction at speed, and becomes both stronger and more efficient above a crossover
  speed. The fluke performs best when its pitch is adapted to the speed so that the angle of attack stays
  near the efficient range of flapping foils (about \SIrange{15}{25}{\degree}).

\item[(RQ3)] \emph{Whole robot.} What cruising speed and power follow for the whole robot?

  \emph{Hypothesis H3.} Both optimized propulsors reach similar cruising speeds, set by the hull and linkage
  resistance, but the fluke needs less power.
\end{description}

RQ1 is answered in \cref{ch:mujoco}. RQ2 and RQ3 are answered in \cref{ch:results}, after the CFD method and
its verification in \cref{ch:cfd,ch:verification}. \Cref{ch:discussion} evaluates the hypotheses.

\section{Contributions}

The contributions of this thesis are:
\begin{enumerate}
  \item \emph{A reduced-order gait-optimization pipeline with an actuator model} (\cref{ch:mujoco}).
        Morison-type link forces with an optional flat-plate lift term, a DC-motor torque--speed model of
        the servos and feasibility-first constraints are combined with CMA-ES. In all 101 feasible optima
        obtained with lift in the model and without a limit on the lift share, lift supplies the forward impulse and drag opposes it. A lift-share limit
        allows drag-based and lift-based strokes to be compared within the same model.
  \item \emph{An optimized drag-type stroke for the BODY2 fan paddle} (\cref{sec:res-drag}). Overset CFD of the
        real leg with a two-state composite for the folding fan yields a stroke with a trapezoidal velocity
        law, a 12-mm folded width and early closing that produces \SIrange{73}{83}{\milli\newton} per leg
        at rest, 3.1--3.5 times the robot's current gait. The added-mass change that the composite does not resolve
        is bounded analytically with an added mass identified from the CFD.
  \item \emph{A speed-resolved comparison of drag- and lift-type propulsors on the same leg}
        (\cref{sec:res-compare}). Thrust, power and efficiency are computed from rest to
        \SI{0.40}{\metre\per\second}. The drag paddle's speed limit is traced to its recovery, the
        thrust and efficiency crossovers are located at about \SI{0.24}{} and \SI{0.22}{\metre\per\second}, and a speed-scheduled
        fluke pitch raises the fluke efficiency at \SI{0.30}{\metre\per\second} from 0.24 to 0.40.
  \item \emph{Self-propulsion estimates and a fidelity assessment} (\cref{sec:res-selfprop,ch:verification}).
        Cruising speeds are bounded with a range of hull resistances, the CFD set-up is validated against
        a flapping-foil experiment and checked for mesh and cycle convergence, and the thrust deficit of
        the real-time solver Flare Live is quantified on the same fluke motion.
\end{enumerate}

\section{Thesis Outline}

\Cref{ch:related} reviews drag- and lift-based swimming in animals and robots, flapping-foil hydrodynamics
and simulation approaches for swimming legged robots. \Cref{ch:foundations} introduces the quantities and
simple models used throughout: quasi-steady paddling, flapping-foil kinematics and the self-propulsion
balance. \Cref{ch:platform} describes the robot, the two propulsors and their leg motions.
\Cref{ch:mujoco} presents the reduced-order model and the gait optimization. \Cref{ch:cfd} describes the
single-leg CFD and the whole-robot Flare Live model, and \cref{ch:verification} verifies and validates the CFD and
checks the real-time solver against it. \Cref{ch:results} reports the drag-type stroke design, the lift-type fluke,
their comparison, the speed schedule for the fluke pitch and the self-propulsion estimates.
\Cref{ch:discussion} interprets the results and states the limitations, and \cref{ch:conclusion} concludes
and outlines future work. The appendices collect additional data, the coordination results of the reduced-order model
(\cref{app:coord}), the Flare Live results (\cref{app:flare-trends}) and information for reproducing the work.
